\documentclass[a4paper,11pt]{article}
\usepackage[utf8]{inputenc}
\usepackage[french]{babel}
\usepackage[T1]{fontenc}
\usepackage{amsmath,amssymb,amsthm}
\usepackage{geometry}
\geometry{margin=2cm}
\usepackage{enumitem}
\usepackage{hyperref}
\usepackage{booktabs}
\usepackage{array}
\usepackage{xcolor}
\usepackage{listings}
\usepackage{titlesec}
\usepackage{fancyhdr}
\usepackage{lastpage}
\usepackage{graphicx}
\usepackage{etoolbox}
\AtBeginDocument{\shorthandoff{;:!?}}

\titleformat{\section}{\large\bfseries\color{blue!60!black}}{}{0pt}{}[\vspace{-0.5ex}\rule{\textwidth}{0.4pt}]
\titleformat{\subsection}{\bfseries\color{blue!40!black}}{}{0pt}{}
\titleformat{\subsubsection}{\itshape\color{blue!30!black}}{}{0pt}{}

\lstdefinestyle{monstyle}{
  frame=single,
  numbers=left,
  numbersep=5pt,
  breaklines=true,
  basicstyle=\small\ttfamily,
  backgroundcolor=\color{gray!5},
  rulecolor=\color{gray!40},
  columns=fullflexible,
  keepspaces=true,
  showstringspaces=false
}
\lstset{style=monstyle}

\pagestyle{fancy}
\fancyhf{}
\fancyhead[L]{\small STA211 -- R\'eseaux de neurones (cours 1)}
\fancyhead[R]{\small Fiche r\'ecapitulative}
\fancyfoot[C]{\thepage/\pageref{LastPage}}
\renewcommand{\headrulewidth}{0.4pt}

\theoremstyle{definition}
\newtheorem*{definition}{D\'efinition}
\theoremstyle{remark}
\newtheorem*{remark}{Remarque}
\newtheorem*{important}{\`A retenir}

\title{\textbf{Fiche r\'ecapitulative -- STA211}\\[4pt]
\large R\'eseaux de neurones -- Cours 1 \\[2pt]
\small Du neurone formel au perceptron multicouche, r\'etropropagation
et optimisation}
\author{Nicolas Thome (CNAM)}
\date{14 juin 2026}

\begin{document}
\maketitle
\thispagestyle{empty}

\begin{abstract}
Cette fiche pr\'esente les fondements des r\'eseaux de neurones
profonds (Deep Learning), couvrant le neurone formel (McCulloch \&
Pitts, 1943), le perceptron, le perceptron multicouche (MLP),
l'apprentissage par r\'etropropagation du gradient (backpropagation),
les fonctions d'activation (seuil, sigmo\"ide, tanh, softmax),
et les aspects pratiques d'optimisation (descente de gradient
stochastique, r\'egularisation).
\end{abstract}

\vfill
\tableofcontents
\newpage

% ===================================================================
\section{Contexte}
% ===================================================================

Le Deep Learning constitue une rupture dans la reconnaissance des
signaux de bas niveau (images, audio, texte). Avant le DL, les
repr\'esentations interm\'ediaires \'etaient con\c{c}ues
artisanalement (handcrafted features). Depuis le DL, ces
repr\'esentations sont apprises automatiquement, ce qui offre~:
\begin{itemize}
  \item Des performances exp\'erimentales sup\'erieures.
  \item Un apprentissage de repr\'esentations de plus haut niveau
    d'abstraction.
  \item Une m\'ethodologie commune et ind\'ependante du domaine.
\end{itemize}

% ===================================================================
\section{Le neurone formel}
% ===================================================================

\subsection{D\'efinition (McCulloch \& Pitts, 1943)}

Le neurone formel est l'unit\'e de base des r\'eseaux de neurones.
Il r\'ealise une transformation d'une entr\'ee vectorielle
$\mathbf{x} \in \mathbb{R}^m$ en une sortie scalaire
$\hat{y} \in \mathbb{R}$ en deux \'etapes~:

\begin{enumerate}
  \item \textbf{Combinaison lin\'eaire (affine)}~:
    $\displaystyle s = \mathbf{w}^\intercal \mathbf{x} + b =
    \sum_{i=1}^m w_i x_i + b$
  \item \textbf{Activation non lin\'eaire}~:
    $\hat{y} = f(s)$
\end{enumerate}

\noindent o\`u $\mathbf{w} \in \mathbb{R}^m$ est le vecteur de
poids synaptiques, $b \in \mathbb{R}$ le biais, et $f$ la fonction
d'activation.

\subsection{Fonctions d'activation}

\paragraph{Fonction seuil (Heaviside)}~:
\begin{equation}
  H(z) = \begin{cases}
    1 & \text{si } z \geq 0 \\
    0 & \text{sinon}
  \end{cases}
\end{equation}
Interpr\'etation biologique~: le neurone est activ\'e
($\hat{y}=1$) si la somme pond\'er\'ee d\'epasse le seuil $-b$.

\paragraph{Fonction sigmo\"ide}~:
\begin{equation}
  \sigma(z) = \frac{1}{1 + e^{-a z}}
\end{equation}
Propri\'et\'es~: lisse, d\'erivable, interpr\'etation probabiliste
($\hat{y} \sim P(1|\mathbf{x})$). Saturante aux extr\'emit\'es.

\paragraph{Fonction tanh}~:
\begin{equation}
  \tanh(z) = \frac{e^{z} - e^{-z}}{e^{z} + e^{-z}}
\end{equation}
\newpage

\subsection{Classification binaire}

Avec une activation sigmo\"ide, la sortie du neurone
$\hat{y} = \sigma(\mathbf{w}^\intercal \mathbf{x} + b)$ repr\'esente
la probabilit\'e d'appartenance \`a la classe 1~:
\begin{equation}
  \mathbf{x} \text{ est classe 1 si } \mathbf{w}^\intercal
  \mathbf{x} + b > 0 \quad (\hat{y} > 0{,}5)
\end{equation}
La fronti\`ere de d\'ecision est \textbf{lin\'eaire}~:
$\mathbf{w}^\intercal \mathbf{x} + b = 0$.

% ===================================================================
\section{Le perceptron}
% ===================================================================

\subsection{Classification multiclasse}

Pour traiter $K$ classes, on utilise $K$ neurones en sortie
(perceptron). Chaque neurone de sortie $k$ calcule~:
\begin{equation}
  s_k = \mathbf{w}_k^\intercal \mathbf{x} + b_k,
  \quad \hat{y}_k = f(s_k)
\end{equation}

\noindent En \'ecriture matricielle~:
$\mathbf{s} = \mathbf{x} \mathbf{W} + \mathbf{b}$,
o\`u $\mathbf{W} \in \mathbb{R}^{m \times K}$ et
$\mathbf{b} \in \mathbb{R}^{1 \times K}$.

\subsection{Softmax}

Pour la classification multiclasse, on utilise l'activation
\textbf{softmax}~:
\begin{equation}
  \hat{y}_k = \frac{e^{s_k}}{\sum_{k'=1}^K e^{s_{k'}}}
\end{equation}
Interpr\'etation probabiliste~: $\hat{y}_k \sim P(k|\mathbf{x},
\mathbf{W})$. Le perceptron avec softmax \'equivaut \`a la
\textbf{r\'egression logistique} (LR).

\subsection{Limite de la classification lin\'eaire}

Le perceptron monocouche est limit\'e aux fronti\`eres de d\'ecision
lin\'eaires. Le probl\`eme du XOR (ou exclusif) n'est pas
s\'eparable lin\'eairement.

% ===================================================================
\section{Perceptron multicouche (MLP)}
% ===================================================================

\subsection{Architecture}

On ajoute une ou plusieurs \textbf{couches cach\'ees} entre
l'entr\'ee et la sortie~:
\begin{equation}
  \mathbf{h} = f(\mathbf{x} \mathbf{W}_h + \mathbf{b}_h),
  \quad \hat{\mathbf{y}} = f(\mathbf{h} \mathbf{W}_y +
  \mathbf{b}_y)
\end{equation}

\noindent $\mathbf{h} \in \mathbb{R}^L$ est la couche cach\'ee
(projection de $\mathbf{x}$ dans un nouvel espace). Avec une
activation non lin\'eaire $f$, la fronti\`ere de d\'ecision devient
\textbf{non lin\'eaire}.

\subsection{R\'eseaux profonds (DNN)}

En ajoutant davantage de couches cach\'ees, on obtient un r\'eseau
de neurones profond (Deep Neural Network). Chaque couche projette
la couche pr\'ec\'edente dans un nouvel espace, apprenant
progressivement des repr\'esentations interm\'ediaires utiles pour
la t\^ache.

% ===================================================================
\section{Apprentissage}
% ===================================================================

\subsection{Apprentissage supervis\'e}

Soit un ensemble d'apprentissage $\mathcal{A} =
\{(\mathbf{x}_i, \mathbf{y}_i^*)\}_{i=1}^N$, avec
$\mathbf{y}_i^*$ encod\'e en one-hot. On d\'efinit une fonction
de perte $\ell(\hat{\mathbf{y}}_i, \mathbf{y}_i^*)$ et on
cherche $\mathbf{W} = \arg\min \mathcal{L}$.

\subsection{Fonction de perte~: entropie crois\'ee}

Pour la classification, on utilise l'entropie crois\'ee
(Cross-Entropy)~:
\begin{equation}
  \ell_{CE}(\hat{\mathbf{y}}_i, \mathbf{y}_i^*) =
  -\sum_{c=1}^K y_{c,i}^* \log \hat{y}_{c,i} =
  -\log \hat{y}_{c^*, i}
\end{equation}
La perte totale sur l'ensemble d'apprentissage~:
\begin{equation}
  \mathcal{L}_{CE}(\mathbf{W}, \mathbf{b}) =
  \frac{1}{N} \sum_{i=1}^N \ell_{CE}(\hat{\mathbf{y}}_i,
  \mathbf{y}_i^*)
\end{equation}

\subsection{Descente de gradient}

On met \`a jour les param\`etres dans la direction oppos\'ee au
gradient~:
\begin{equation}
  \mathbf{W}^{(t+1)} = \mathbf{W}^{(t)} -
  \eta \frac{\partial \mathcal{L}_{CE}}{\partial \mathbf{W}}
\end{equation}
o\`u $\eta$ est le taux d'apprentissage (learning rate).

\subsection{R\'etropropagation (backpropagation)}

Le gradient se calcule par \textbf{r\'etropropagation} de l'erreur
en utilisant la r\`egle de d\'erivation en cha\^ine~:
\begin{equation}
  \frac{\partial \ell}{\partial x} =
  \frac{\partial \ell}{\partial \hat{y}}
  \frac{\partial \hat{y}}{\partial x}
\end{equation}

\paragraph{Pour la r\'egression logistique}~:
\begin{align}
  \frac{\partial \ell_{CE}}{\partial \hat{y}_i} &=
  -\frac{1}{\hat{y}_{c^*, i}} \\
  \frac{\partial \ell_{CE}}{\partial \mathbf{s}_i} &= \hat{\mathbf{y}}_i -
  \mathbf{y}_i^* = \boldsymbol\delta_i^y \\
  \frac{\partial \ell_{CE}}{\partial \mathbf{W}} &=
  \mathbf{x}_i^\intercal \boldsymbol\delta_i^y
\end{align}

\paragraph{Pour le perceptron (1 couche cach\'ee)}~:
\begin{align}
  \boldsymbol\delta_i^h &= \boldsymbol\delta_i^y \mathbf{W}_y^\intercal
  \odot \mathbf{h}_i \odot (1 - \mathbf{h}_i) \\
  \frac{\partial \ell}{\partial \mathbf{W}_h} &=
  \mathbf{x}_i^\intercal \boldsymbol\delta_i^h \\
  \frac{\partial \ell}{\partial \mathbf{W}_y} &=
  \mathbf{h}_i^\intercal \boldsymbol\delta_i^y
\end{align}

\paragraph{Cas g\'en\'eral (couche $l$)}~:
\begin{align}
  \frac{\partial \mathcal{L}}{\partial \mathbf{U}_l} &=
  \boldsymbol\Delta_l = \boldsymbol\Delta_{l+1} \mathbf{W}_{l+1}^\intercal
  \odot \mathbf{H}_l \odot (1 - \mathbf{H}_l) \\
  \frac{\partial \mathcal{L}}{\partial \mathbf{W}_l} &=
  \mathbf{H}_{l-1}^\intercal \boldsymbol\Delta_l
\end{align}

\subsection{Descente de gradient stochastique (SGD)}

Le calcul du gradient complet $\nabla_{\mathbf{w}}$ sur tout
l'ensemble d'apprentissage est trop co\^uteux. Approximations~:
\begin{itemize}
  \item \textbf{En ligne (online)}~: un seul exemple par mise \`a jour.
  \item \textbf{Minibatch}~: $B < N$ exemples par mise \`a jour.
\end{itemize}
Le gradient bruit\'e permet plus de mises \`a jour et une
convergence plus rapide en pratique.

\subsection{Plan de taux d'apprentissage (learning rate decay)}

Le taux d'apprentissage $\eta$ peut d\'ecro\^itre au cours de
l'entra\^inement~:
\begin{itemize}
  \item D\'ecroissance inverse~: $\eta_t = \frac{\eta_0}{1 + r \cdot t}$
  \item D\'ecroissance exponentielle~: $\eta_t = \eta_0 e^{-\lambda t}$
  \item D\'ecroissance par paliers~: $\eta_t = \eta_0 r^{\lfloor t / t_u \rfloor}$
\end{itemize}

% ===================================================================
\section{R\'egularisation et g\'en\'eralisation}
% ===================================================================

\subsection{Sur-apprentissage (overfitting)}
L'objectif n'est pas de minimiser la perte sur l'ensemble
d'apprentissage, mais de g\'en\'eraliser \`a des donn\'ees non
vues. Un \'ecart \'elev\'e entre la perte d'apprentissage et la
perte de test indique du sur-apprentissage.

\subsection{R\'egularisation}

Ajout d'un terme de p\'enalit\'e $R(\mathbf{W})$ dans l'objectif~:
\begin{equation}
  \mathcal{L}(\mathbf{W}) = \mathcal{L}_{CE}(\mathbf{W}) +
  \alpha R(\mathbf{W})
\end{equation}

\begin{itemize}
  \item \textbf{$L_2$ (weight decay)}~: $R(\mathbf{W}) =
    \|\mathbf{W}\|^2$ (le plus courant).
  \item \textbf{$L_1$}~: $R(\mathbf{W}) = \|\mathbf{W}\|_1$
    (favorise la parcimonie).
  \item \textbf{Dropout}~: d\'esactive al\'eatoirement des neurones
    pendant l'entra\^inement.
\end{itemize}

\subsection{Hyper-param\`etres}

R\'eseaux de neurones~: nombreux hyper-param\`etres \`a r\'egler~:
\begin{itemize}
  \item Param\`etres d'optimisation~: learning rate, weight decay,
    plan de decay, nombre d'\'epoques.
  \item Param\`etres d'architecture~: nombre de couches, nombre de
    neurones, type de non-lin\'earit\'e.
\end{itemize}
L'ajustement se fait sur un ensemble de validation.

% ===================================================================
\newpage
\section*{Questions cl\'es (QCM)}
% ===================================================================
\addcontentsline{toc}{section}{Questions cl\'es (QCM)}

\begin{enumerate}

  \item Quelle est la sortie d'un neurone formel~?
    \begin{enumerate}
      \item $\hat{y} = \mathbf{w}^\intercal \mathbf{x} + b$
      \item $\hat{y} = f(\mathbf{w}^\intercal \mathbf{x} + b)$
      \item $\hat{y} = \mathbf{x} \mathbf{W} + \mathbf{b}$
      \item $\hat{y} = f(\mathbf{x})$
    \end{enumerate}
    \textbf{R\'eponse~: b (combinaison lin\'eaire puis activation).}

  \item Quelle fonction d'activation donne une interpr\'etation
    probabiliste en classification binaire~?
    \begin{enumerate}
      \item Seuil (Heaviside)
      \item Sigmo\"ide
      \item Tanh
      \item ReLU
    \end{enumerate}
    \textbf{R\'eponse~: b (sigmo\"ide, $\hat{y} \sim P(1|\mathbf{x})$).}

  \item Quelle activation est utilis\'ee pour la classification
    multiclasse~?
    \begin{enumerate}
      \item Sigmo\"ide
      \item Tanh
      \item Softmax
      \item Seuil
    \end{enumerate}
    \textbf{R\'eponse~: c (softmax).}

  \item Le perceptron monocouche \'equivaut \`a~:
    \begin{enumerate}
      \item La r\'egression lin\'eaire
      \item La r\'egression logistique
      \item L'analyse discriminante
      \item Les SVMs
    \end{enumerate}
    \textbf{R\'eponse~: b (r\'egression logistique).}

  \item Quel probl\`eme illustre la limite de la classification
    lin\'eaire~?
    \begin{enumerate}
      \item Le XOR
      \item Le AND
      \item Le OR
      \item Le clustering
    \end{enumerate}
    \textbf{R\'eponse~: a (XOR, non s\'eparable lin\'eairement).}

  \item Comment s'appelle la technique de calcul du gradient dans
    un MLP~?
    \begin{enumerate}
      \item Gradient boosting
      \item R\'etropropagation (backpropagation)
      \item Descente de gradient
      \item Forward propagation
    \end{enumerate}
    \textbf{R\'eponse~: b (r\'etropropagation du gradient).}

  \item Quelle approximation du gradient est utilis\'ee pour les
    grands jeux de donn\'ees~?
    \begin{enumerate}
      \item Gradient complet
      \item SGD (minibatch)
      \item Newton-Raphson
      \item BFGS
    \end{enumerate}
    \textbf{R\'eponse~: b (SGD minibatch).}

  \item Que p\'enalise la r\'egularisation $L_2$ (weight decay)~?
    \begin{enumerate}
      \item La somme des poids en valeur absolue
      \item Le carr\'e de la norme des poids
      \item Le nombre de couches
      \item Le taux d'apprentissage
    \end{enumerate}
    \textbf{R\'eponse~: b (carr\'e de la norme $\|\mathbf{W}\|^2$).}

  \item Quel est le r\^ole d'un ensemble de validation~?
    \begin{enumerate}
      \item Apprendre les param\`etres
      \item R\'egler les hyper-param\`etres
      \item Calculer la perte finale
      \item Initialiser les poids
    \end{enumerate}
    \textbf{R\'eponse~: b (r\'egler les hyper-param\`etres).}

  \item Qui a propos\'e le neurone formel en 1943~?
    \begin{enumerate}
      \item Rosenblatt
      \item McCulloch \& Pitts
      \item LeCun
      \item Hinton
    \end{enumerate}
    \textbf{R\'eponse~: b (McCulloch \& Pitts, 1943).}

\end{enumerate}

% ===================================================================
\begin{thebibliography}{9}
% ===================================================================

\bibitem{MP43}
  W.~S.~McCulloch, W.~Pitts.
  \newblock \emph{A logical calculus of the ideas immanent in
    nervous activity}.
  \newblock Bull. Math. Biophysics, 5:115--133, 1943.

\bibitem{Rosenblatt1958}
  F.~Rosenblatt.
  \newblock \emph{The perceptron: A probabilistic model for
    information storage and organization in the brain}.
  \newblock Psychological Review, 65(6):386--408, 1958.

\bibitem{Werbos1974}
  P.~Werbos.
  \newblock \emph{Beyond regression: New tools for prediction and
    analysis in the behavioral sciences}.
  \newblock Ph.D. thesis, Harvard University, 1974.

\bibitem{Cybenko1989}
  G.~Cybenko.
  \newblock \emph{Approximation by superpositions of a sigmoidal
    function}.
  \newblock Math. Control Signals Systems, 2(4):303--314, 1989.

\bibitem{LeCun1989}
  Y.~LeCun et al.
  \newblock \emph{Backpropagation applied to handwritten zip code
    recognition}.
  \newblock Neural Computation, 1(4):541--551, 1989.

\end{thebibliography}

\end{document}
